\section{Division-free circuits and arithmetic savings}
\label{sec:circuits}

All displayed auxiliary symbols in this section denote polynomial
abbreviations. They introduce neither existentially quantified coordinates
nor operations outside addition, subtraction, and multiplication.
The delivered JSON files specify their evaluation order and every operand.
The accounting convention charges each binary arithmetic operation, including
multiplication by a fixed integer; reusing an earlier value and comparing two
computed expressions are free. This is the convention used by the repository
comparison circuits \cite{repo21,repo22}.
All polynomial-identity comparisons hold the four coefficient ports equal.
Replacing the whole tuple is a separate semantic step proved in
Section~\ref{sec:coefficients}.

\subsection{A shared evaluator for the four-root construction}

Use the notation of the root-recovery construction, and write $W$ for its
supplied merged expression, $N$ for the weight of the main root, and $m=HM$
for the weight of the first Pell root. Thus the relevant eight-conjugate
product is
\begin{equation}
\mathcal R=
\prod_{\sigma,\epsilon,\delta\in\{-1,1\}}
\left[d\bigl(W+\sigma N\sqrt G+\epsilon m\sqrt A\bigr)^2
      -H^2e+\delta H^2B_{\rm aux}\sqrt D\right].
\label{eq:circuits-eight-norm}
\end{equation}
Here $B_{\rm aux}=2Qvb$. Define $P_a=QDv^2$ and $U_a=Qb^2$,
so that $B_{\rm aux}^2D=4P_aU_a$. Formula \eqref{eq:circuits-eight-norm} explains the elimination;
it is not the circuit definition and does not require a square-root operation.

The following completely integral evaluator computes the same polynomial:
\begin{align*}
 V_0&=W^2+N^2G,& J_\theta&=4W^2(N^2G),& a_0&=m^2A,\\
 C_0&=d(V_0+a_0)-H^2e,& K_*&=4d^2a_0,& U_0&=K_*V_0,\\
 Z_0&=4H^4P_aU_a,\\
 x&=C_0^2+J_\theta d^2-U_0-Z_0,
 &y&=2C_0d-K_*,\\
 n_0&=x^2+J_\theta y^2-4Z_0U_0,
 &n_1&=2xy-4Z_0K_*,\\
 \mathcal R&=n_0^2-J_\theta n_1^2.
\end{align*}
The three-equation system compares
\[
 N_k=1,\qquad N_t=1,\qquad n_0^2=J_\theta n_1^2.
\]
The corresponding single equation is
\[
 N_kN_t(1+\mathcal R)-1=0.
\]
Since $4\mid J_\theta$, the integer $\mathcal R$ is a square modulo $4$.
Consequently $1+\mathcal R$ cannot equal $-1$. This preserves the modular
guard needed in the product equation; the remaining outer signs are handled
by the root-recovery proof.

To verify the evaluator, work temporarily in a quadratic algebra and put
$\theta=2WN\sqrt G$, so that $\theta^2=J_\theta$ and
$(W+N\sqrt G)^2=V_0+\theta$. Taking the norm in the first-root direction
uses the exact identity
\begin{equation}
 \bigl((C-u)^2-Z\bigr)\bigl((C+u)^2-Z\bigr)
   =(C^2-u^2-Z)^2-4u^2Z.
\label{eq:circuits-symmetric}
\end{equation}
After the first two eliminations, the remaining quadratic-algebra element
is $n_0+n_1\theta$. Its norm is the last line of the evaluator. This is an
identity in a free integer-polynomial ring. In particular no division by
$W$, $N$, or $2WN$ is performed, and the identity remains valid when one of
these quantities is zero.

\subsection{Factoring the common odd coefficient}

The earlier paired evaluator explicitly constructs $V_1=2WN$, then
$dV_1$, $K_*V_1$, and their products and squares. Every odd component has
this same factor $V_1$. Retaining only its square
$G V_1^2=4W^2(N^2G)$ avoids three arithmetic operations in the complete
polynomial evaluator. Both $W^2$ and $N^2G$ are already needed for $V_0$.
The system saves one further subtraction by comparing its two final norm
terms directly.

Applied to the unchanged nested-norm polynomial of package~21, this gives
\[
 (12\hbox{ positive variables},\ 1\hbox{ equation},\ 1353,\ 113)
\]
in the order (unknowns, equations, degree, operations), improving its
116-operation evaluator. Its three-equation system uses 108 operations
rather than 112, with the same maximum degree 1344. The polynomial itself,
its zero set, and the proof assumptions are unchanged \cite{repo21}.

The same idea applies to the monic-resultant circuits of package~22.
There $r_1=2W$ occurs in every odd component, while $Br_1^2$ already occurs
in the even component. In its twelve-variable evaluator, write
\[
 v_I=r_1t,\qquad
 h_I=2r_1(v_Rt-\lambda u_R).
\]
One may compute $Bv_I^2=(Br_1^2)t^2$ and
$Bh_I^2=4(Br_1^2)(v_Rt-\lambda u_R)^2$ directly. Moreover
$Br_1^2=4BW^2$ needs two gates because $W^2$ is already available,
whereas forming $r_1$, squaring it, and multiplying by $B$ uses three.
The complete saving is three operations: 117 becomes 114 for the
polynomial, and 112 becomes 109 for its system. The thirteen-variable
intermediate has $h_I=r_1t$ and saves two operations: its degree-319
polynomial uses 89 operations, and its four-equation system of maximum
degree 288 uses 81. Again these are exact refactorings of the original
polynomials \cite{repo22}.

\subsection{Weights and the unrestricted-integer lift}

The new positive constructions use the same shared evaluator with different
separating weights. For the degree-1081 version, the identity
\[
 1+8TI+16REi^2c^3=1+8I(T+2REic),\qquad I=ic^2,
\]
saves one more operation. The tighter degree-1073 and degree-1069 weights
are evaluated in their already factored form
\[
 M=I(2T+R\mathcal C\,ic),
\]
with $\mathcal C=E+3(X+Y)$ and
$\mathcal C=E+2(X+Y)+sq^2$, respectively. Their improvement in degree
comes from the cancellation proved in the weight analysis, rather than
from the generic degree bound of the evaluation graph.

For the signed construction use
\[
 M=1+8(TI)^2+16Ra i^2c^3
   =1+8(ic)I(T^2c+2Ra),\qquad N=1+8HM.
\]
The second expression is a polynomial identity for arbitrary integers,
including negative $i$ and $c$. The separate signed-domain proof supplies
the positivity properties needed in the root argument.

There is a further saving when the six guarded coordinates are replaced by
$\Pos(u,v,t)=1+u^2+v^2+t(t+1)$. Write
\[
 F=F_0+1,\quad z=z_0+1,\quad q=\beta J+1,
\]
where $F_0$ and $z_0$ are their three-term nonnegative quadratic parts.
Then
\begin{align*}
 q-F&=\beta J-F_0,\\
 (K_0+w)C+(q-F)-z(q-1)
   &=(K_0+w)C-F_0-z_0\beta J.
\end{align*}
Thus the evaluator never forms the two separate $+1$ gates for $F$ and
$z$. This changes neither the substitution nor the quantified domain.
The final unrestricted-integer polynomial uses 167 operations, and its
three-equation system uses 162.

\subsection{Final ledgers and verification scope}

\begin{center}
\small
\begin{tabular}{lrrrr}
\hline
Form & Polynomial degree & Operations & System degree & Operations\\
\hline
12 positive, tightest weight &1069&123&1060&118\\
12 positive, tight weight    &1073&121&1064&116\\
12 positive, simpler weight &1081&120&1072&115\\
12 positive, refactored norm&1353&113&1344&108\\
13 positive, monic merger   &319&89&288&81\\
24 unrestricted integers    &1936&167&1920&162\\
\hline
\end{tabular}
\end{center}
All systems in this table have three equations except the thirteen-variable
system, which has four. The degrees are exact. For the degree-1069 and
1073 products, naive propagation through the graph gives 1081; their
homogeneous cancellation must be used to obtain the exact degree.
The graph metadata records this distinction rather than overwriting the
larger syntactic bound.

The package contains three complementary kinds of evidence. First,
\path{code/check_exact_identities.py} expands the small free-polynomial
identities over $\Z$ and checks every coefficient. This covers the shared
odd-component factors, identity~\eqref{eq:circuits-symmetric}, the weight
factorizations, and both integer-lift cancellations. Second, independent
exact evaluations multiply all eight conjugates in an eight-dimensional
algebra and compare that result with the division-free graph. Third,
\texttt{degree\_check.py} evaluates the entire graph on univariate slices
over two prime fields. Its convolution uses integer packing with a proved
carry bound; it does not use floating point or merely propagate the
largest input degree. These checks confirm the degree cancellations and
support the separate uniform leading-form proofs.

For the signed polynomial, the independently derived monomial
\[
 J_x^{1354}h\,i_{\rm int}^{48}\eta_x^{192}
 w_{\rm int}^{162}s_{\rm int}^{177}z_x^2
\]
has total degree 1936 and coefficient
\[
 2^{257}\,255^4(57599\cdot73983)^2\beta^{677}.
\]
The coefficient is nonzero for every admissible $\beta>0$; its modular
values agree with two independent full-polynomial slice calculations.

Arithmetic identity checks establish the stated circuit transformations.
The root-recovery lemmas, the signed-domain argument, the native compiler,
and the calibration of K215 remain mathematical proof dependencies. None
of these finite calculations is presented as a new Lean or Rocq theorem,
or as a test of whether the consistency polynomial has a solution.
